%% For tips on how to write a great abstract, have a look at
%%	-	https://www.cdc.gov/stdconference/2018/How-to-Write-an-Abstract_v4.pdf (presentation, start here)
%%	-	https://users.ece.cmu.edu/~koopman/essays/abstract.html
%%	-	https://search.proquest.com/docview/1417403858
%%  - 	https://www.sciencedirect.com/science/article/pii/S037837821830402X

\begin{abstract}
Instructional elements in Christian digital games are a frequent source of complaint among players. Nevertheless, the reason for such critiques remains unknown. Related research either avoids the topic entirely, mentions it without investigating its causes, or examines the success of Christian digital games more generally. Those in the last category mostly list factors outside developers’ control and fail to adequately consider the games’ content. 

My approach combines theories of didacticism, heavy-handedness, and player involvement to create a framework that offers two complementary explanations of how a Christian digital game’s informative formal elements might lead to negative player experiences. The first contends that overly instructive content might lead players to presume that the beliefs being presented are distorted by the developer’s cognitive bias. The second notes that designers can implement instruction in ways that limit involving affordances.

Using this framework, I perform a game analysis of two Christian digital games, \textit{Alum} and \textit{Captain Bible in Dome of Darkness}, which I supplement with public player reviews. For \textit{Alum}, I observe that the game contains multiple scripted sequences that give players grounds on which to believe the author is a prejudiced zealot. Meanwhile, though some players complain about the game’s overtness, none of these complaints can be explained exclusively by the game restricting players’ interpretative agency. On the other hand, the pedagogical gameplay loop in \textit{Captain Bible}, along with several other formal elements, seems to restrict players’ ludic and affective involvement. This caused some players to feel as if the gameplay was too monotonous and the levels too barren.

Given the lack of examples, I avoided generalising my observations.  Nevertheless, the framework seems to give clear reasons for players' grievances with edifying content in these examples. Thus, their designers can attempt to resolve such problems with more confidence. Lastly, I noted that while some solutions might be constrained by external factors such as finances, others are not. Therefore, I suggest that the framework presented in this study may be useful not only for Christian digital game developers, but also for analysing other games plagued by similar criticism.

\end{abstract}

